\chapter{Introdução}
\label{cap:introducao}

A preparação para concursos públicos caracteriza-se como um processo educacional intensivo e de alto risco, no qual a simples dedicação de tempo não garante o sucesso. A jornada do estudante é marcada pela necessidade de dominar grandes volumes de conteúdo e de converter esse conhecimento em resultados sob a pressão do exame. Portanto, o tema deste trabalho é o desenvolvimento de uma solução tecnológica que aborda a gestão do processo de estudo. A eficácia na preparação exige que o estudante demonstre tanto competência no material quanto a habilidade de lidar com o contexto de avaliação. A Competência em Testes e a Gestão do Tempo estão entre os fatores com as correlações mais significativas com a performance acadêmica \cite{kleijn_1994}, evidenciando que a falha em qualquer um destes componentes é uma causa primária de baixo rendimento.

As causas do baixo desempenho na preparação podem ser rastreadas em hábitos de estudo ineficazes. Muitos estudantes negligenciam o valor da prática de recuperação ativa e da autoavaliação contínua, utilizando estratégias passivas que não se traduzem em sucesso na hora da prova. O desenvolvimento de um alto desempenho está ligado à execução de estratégias de estudo específicas. Por exemplo, estudos que utilizam o inventário LASSI (\textit{Learning and Study Strategies Inventory}) apontaram que as subescalas de Estratégias de Teste e Concentração são significativos para a performance acadêmica \cite{zhou_2016}. A ausência de um sistema estruturado para fomentar o \textit{self-testing} (como flashcards inteligentes) e a falta de foco (como o método Pomodoro) são, portanto, causas comportamentais diretas da ineficiência do estudo.

Como consequência direta do estudo desorganizado e da pressão proveniente dos concursos, os estudantes enfrentam uma elevada pressão psicológica. A gestão ineficiente do tempo e a baixa confiança nas estratégias de estudo resultam em altos níveis de ansiedade e estresse. Este quadro psicológico negativo tem um efeito prejudicial no desempenho. Dessa maneira, a autorregulação e a resiliência são destacadas como dimensões cruciais que mitigam esses desafios afetivos, sugerindo que a preparação não pode ignorar a saúde mental do candidato \cite{liu_2025}.

Diante desses desafios, o mercado de aplicativos de estudo tem crescido, mas geralmente de forma segmentada. Existem aplicativos dedicados à gestão do tempo (Pomodoro), outros à memorização (Flashcards) e outros à motivação (Gamificação). Contudo, essa fragmentação exige que o usuário integre manualmente diferentes ferramentas, limitando a eficácia e a análise geral do progresso. A oportunidade de inovação surge, portanto, na unificação inteligente dessas funcionalidades. A necessidade de incluir a Gamificação é corroborada pela literatura, que reconhece o valor da integração da mesma para a otimização do e-learning e o aumento do engajamento \cite{marengo_2025}.

Portanto, este trabalho se posiciona para preencher essa lacuna ao desenvolver um aplicativo integrado e inteligente. Dessa maneira, o trabalho propõe o Desenvolvimento de um Aplicativo Mobile Gratuito para Otimização do Aprendizado em Preparação para Concursos Públicos, com o objetivo de integrar as metodologias de estudo mais validadas pela pesquisa acadêmica em uma única aplicação gratuita. As estratégias de aprendizado e os hábitos de estudo estão diretamente ligados ao desempenho acadêmico \cite{bickerdike_2016}. O presente trabalho busca, assim, desenvolver a engenharia de software para uma ferramenta que atenda à demanda por uma preparação eficaz, organizada e psicologicamente sustentável.

\section{Contextualização}

A preparação para concursos públicos ou provas insere-se no campo da educação de pessoas, onde a excelência no desempenho é uma necessidade. O estudante, muitas vezes inserido em uma rotina de múltiplas demandas, precisa converter tempo de estudo em retenção de longo prazo e sucesso nessas avaliações. Este cenário competitivo exige que o candidato desenvolva, de forma consciente, o que a literatura denomina como estratégias de aprendizagem e hábitos de estudo. A relevância de uma abordagem focada em habilidades, e não apenas em conteúdo, é suportada por estudos que confirmam uma associação significativa entre as estratégias de aprendizado, os hábitos de estudo e o desempenho acadêmico \cite{bickerdike_2016}. Portanto, o desenvolvimento de uma ferramenta digital deve ter como foco principal a otimização dessas metodologias.

Um dos principais desafios que a preparação prolongada impõe é a gestão da ansiedade e o foco. O estudo desorganizado e a sensação de perda de controle sobre o volume de material são fatores que intensificam a tensão. Embora o efeito prejudicial da ansiedade no desempenho seja um fenômeno documentado, a literatura sugere que para melhorar o desempenho deve-se focar nas causas subjacentes da ansiedade, que geralmente estão ligadas à falta de técnicas de estudo adequadas \cite{kleijn_1994}. Isso justifica a necessidade de o aplicativo integrar recursos de foco, como o Pomodoro, e de organização de tarefas, visando sanar as causas comportamentais da ansiedade, e não apenas seus sintomas.

No domínio cognitivo, o sucesso na prova não advém da quantidade de informações absorvidas, mas sim da eficiência em recuperá-las eficazmente. A eficácia das técnicas de estudo integradas propostas no aplicativo, como Flashcards, é comprovada pela sua capacidade de forçar a recuperação ativa. Ao analisar os fatores de desempenho, foi demonstrado que as subescalas de Gerenciamento do Tempo (\textit{Time Management}) e Autoavaliação (\textit{Self-Testing}) são fatores mais fortes no desempenho acadêmico \cite{zhou_2016}. Isso valida o foco em funcionalidades que promovam a autodisciplina e a prática constante de recuperação de informações, elementos centrais na autorregulação e na retenção.

Finalmente, a sustentabilidade da preparação para provas depende da atenção ao Bem-Estar e à saúde mental. O estudante não pode ser tratado apenas como um recipiente de informações, mas como um indivíduo que precisa de suporte para lidar com o estresse. Os recursos de check-ins de humor e dicas psicológicas se baseiam na ideia de que a resiliência é um fator protetor. Estudos sobre essa área indicam que as três dimensões da resiliência na aprendizagem (autorregulação, sociabilidade e empatia) têm um efeito positivo direto no desempenho acadêmico \cite{liu_2025}. Ao fornecer ferramentas para que o estudante monitore e gerencie sua autorregulação e seu estado emocional, o aplicativo se estabelece como um facilitador de um estudo mais equilibrado e, consequentemente, mais eficiente.

\section{Justificativa}

O desenvolvimento deste aplicativo se justifica primariamente pela necessidade de superar a ineficiência inerente aos métodos de estudo passivos e desorganizados, que são frequentemente adotados pelos estudantes. A mera dedicação de horas à leitura de materiais não garante a retenção nem a aplicação do conhecimento em situações de teste. O trabalho se propõe a transformar o estudo em um processo ativo e estruturado, baseado em evidências de sucesso. A urgência dessa intervenção é reforçada pela pesquisa, que demonstra que o aprimoramento das habilidades de estratégias de estudo/aprendizagem está associado ao sucesso acadêmico tanto em exames internos quanto em exames externos \cite{khalil_2024}.

O principal foco que o projeto visa englobar é a falta de um sistema integrado que promova a autorregulação completa do estudante. O sucesso acadêmico depende de três componentes estratégicos: \textit{Skill} (habilidade), \textit{Will} (vontade) e \textit{Self-Regulation} (autorregulação). As funcionalidades de Gestão de Cronograma, Pomodoro e \textit{Self-Testing} (Autoavaliação) do aplicativo visam construir esse último componente. A literatura, ao analisar o Inventário de Estratégias de Aprendizagem e Estudo (LASSI), especifica que as subescalas de Gerenciamento do Tempo (\textit{Time Management}) e Autoavaliação (\textit{Self-Testing}) são fatores mais fortes no desempenho acadêmico em períodos iniciais \cite{zhou_2016}. Isso justifica a seleção destas funcionalidades como elementos centrais e não periféricos do software.

No mercado atual de e-learning, as ferramentas de otimização de estudo são predominantemente fragmentadas. Existem aplicativos de pomodoro, gerenciadores de tarefas, e plataformas de flashcards, mas a ausência de uma plataforma unificada que colete dados de todas essas interações (tempo, acertos, motivação) impede uma visão holística e adaptativa do progresso. A justificativa técnica e de inovação deste trabalho reside exatamente na unificação dessas áreas com a preocupação psicológica do estudante, contendo check-ins de humor e dicas de bem-estar para regular o nível de ansiedade e estresse do mesmo. Outra adição do aplicativo é a gamificação, que é uma estratégia validada para a otimização do e-learning e aumento do engajamento \cite{marengo_2025}.

Finalmente, o desenvolvimento deste trabalho representa uma contribuição prática e teórica significativa para a área de Engenharia de Software e Educação. Ao projetar uma engenharia de software para um aplicativo gratuito e baseado em pesquisa, o trabalho democratiza o acesso a metodologias de estudo de alta eficácia. Este projeto é, portanto, uma aplicação tangível e socialmente relevante da ciência da computação para a otimização educacional.

\section{Questão de Pesquisa}

Para fundamentar o desenvolvimento deste trabalho, a questão de pesquisa foi formalmente definida: Como promover intervenções personalizadas na rotina de estudos de estudantes por meio de um aplicativo mobile gratuito baseado em metodologias eficazes, gamificação e controle psicológico?

\section{Objetivos}

Este trabalho possui um objetivo geral e outros seis objetivos específicos que serão aprofundados nesta seção.

\subsection{Objetivo Geral}

Este trabalho possui como objetivo geral o desenvolvimento de um aplicativo \textit{open source} mobile gratuito que visa o aprimoramento da rotina de estudos de estudantes concurseiros, por meio de metodologias eficazes, gamificação e controle psicológico para o bem-estar.

\subsection{Objetivos Específicos}

São objetivos específicos deste trabalho:

\begin{itemize}

\item Fornecer metodologias de estudo eficazes e comprovadas pela pesquisa aos estudantes;

\item Aplicar o conceito de gamificação no aplicativo a fim de tornar a jornada mais motivadora;

\item Oferecer orientações personalizadas ao notar níveis de estresse elevados;

\item Disponibilizar gráficos de progressos individuais para os usuários do aplicativo;

\item Integrar todos os objetivos acima em um único aplicativo gratuito;

\item Servir como uma ferramenta gratuita para disseminação de bons hábitos de aprendizado.

\end{itemize}

\section{Organização do Trabalho}

O presente trabalho está dividido nos seguintes capítulos:

\begin{itemize}

\item \textbf{Capítulo 1 - Introdução:} A função do capítulo introdutório é preparar o leitor ao descrever o cenário em que a proposta será implementada, apresentar o porquê da pesquisa, mostrar a pergunta-chave do trabalho, listar os resultados esperados e resumir as estratégias empregadas.

\item \textbf{Capítulo 2 - Referencial Teórico:} Dividido em pontos-chave, este capítulo desenvolve a base teórica do estudo, introduzindo conceitos importantes relacionados à temática. Os tópicos abordados englobam desde a gestão de tempo e metodologias de estudo até a gamificação, além de uma revisão sobre aplicativos já existentes e noções afins.

\item \textbf{Capítulo 3 - Metodologia:} Neste capítulo, é detalhado a metodologia de pesquisa e desenvolvimento do projeto. Nele, são explicadas as abordagens que garantiram a implementação eficaz da ferramenta de software.

\item \textbf{Capítulo 4 - Desenvolvimento:} O capítulo aprofunda os detalhes da fase de desenvolvimento do software, incluindo as otimizações e modificações realizadas para assegurar uma implementação de sucesso.

\item \textbf{Capítulo 5 - Considerações Finais:} O capítulo final resume os resultados obtidos pelo trabalho e indica sugestões para projetos ou estudos subsequentes.

\end{itemize}
